\section{边界情形与证明核查}
新增检查包括完整 ResNet-18 的参数量与八个残差块、卷积参数组的无重复覆盖、$K=0$ 的准确含义和零梯度行为、谱范数误差下降不等式的数值核查、两因子目标的联合有限差分梯度，以及数据增强的种子复现。它们补充下文原有版本的检查项目。
\subsection{新增实验的数值汇总}
\begin{table}[htbp]\centering\small
\caption{新增病态矩阵分解的末点指标，三个种子的均值与样本标准差。}
\begin{tabular}{lcc}\toprule
方法 & 末点目标值 & 相对重构误差 \\ \midrule
SGD + momentum & $1.83\times10^{-7}\,\pm\,1.1\times10^{-7}$ & $3.91\times10^{-4}\,\pm\,1.3\times10^{-4}$ \\
AdamW & $2.89\times10^{-6}\,\pm\,2.4\times10^{-6}$ & $1.52\times10^{-3}\,\pm\,6.4\times10^{-4}$ \\
Muon (NS5) & $1.02\times10^{-4}\,\pm\,2.7\times10^{-5}$ & $9.49\times10^{-3}\,\pm\,1.3\times10^{-3}$ \\
\bottomrule\end{tabular}

\end{table}
\begin{table}[htbp]\centering\small
\caption{ResNet-18 的 $K$ 消融数值。学习率固定为 $K=5$ 试跑选定值，其余训练设置完全相同。}
\begin{tabular}{cccc}\toprule
$K$ & 末轮验证 Acc / \% & 选定测试 Acc / \% & ms/step \\ \midrule
0 & $69.440\,\pm\,0.801$ & $68.863\,\pm\,0.811$ & $10.166\,\pm\,0.083$ \\
1 & $68.927\,\pm\,0.814$ & $68.453\,\pm\,0.186$ & $13.484\,\pm\,0.047$ \\
3 & $75.580\,\pm\,0.208$ & $74.507\,\pm\,0.430$ & $20.316\,\pm\,0.104$ \\
5 & $75.420\,\pm\,0.517$ & $74.567\,\pm\,0.271$ & $26.846\,\pm\,0.216$ \\
7 & $75.633\,\pm\,0.244$ & $74.527\,\pm\,0.485$ & $33.671\,\pm\,0.059$ \\
\bottomrule\end{tabular}

\end{table}
\subsection{一个不应使用的矩阵范数不等式}
对可逆对角矩阵 $S=\operatorname{diag}(1,2)$，
\begin{equation}
 \op{S^{-1}}=1,\qquad \frac1{\op S}=\frac12.
\end{equation}
因此一般不能使用 $\|S^{-1}\|_{\mathrm{op}}\leq1/\|S\|_{\mathrm{op}}$。正确关系是 $\|S^{-1}\|_{\mathrm{op}}=1/s_{\min}(S)$。本文用式\eqref{eq:alignment}替代含逆矩阵的论证，同时避免要求梯度或动量始终满秩。指出这一中间步骤的问题，不等于断言所有相关定理都无法通过其他方法证明。

\subsection{数值秩与浮点误差}
精确数学定义只保留正奇异值。程序采用阈值
\begin{equation}
 s_i>\max(m,n)\,\epsilon_{\mathrm{mach}}\,s_1,
\end{equation}
以避免浮点 SVD 把数值零误认为有效方向。因此实现严格说是数值秩上的极因子。合成实验使用 float64；语言模型的 SVD 对照使用 float32。阈值附近的方向可能与实数模型不同，这是实现限制，不能通过把容差设置为零来消除。

\subsection{检查项目与实验解释}
单元检查包括：零秩、秩一与满秩矩阵上的对偶恒等式和范数界；二次目标的有限差分梯度；带任意跟踪误差的单步下降；有限多项式与标量奇异值迭代的一致性；更新时不原地改写梯度；参数分组无重复且覆盖全部参数；因果注意力不使用未来字符；同一种子训练一步的精确复现。

有限步定理控制的是理论期望，随机实验只提供有限样本观察。随机实验的 CSV 内包含真实梯度，是因为合成目标允许计算该量；这并不意味着实际大模型训练时已知全数据梯度。语言模型没有验证全局光滑性或全程方差界，故其曲线只能作为应用实验报告。

\section{运行配置与文件对应}
\begin{table}[htbp]\centering\small
\caption{主要运行记录及其用途}
\begin{tabularx}{\linewidth}{lX}\toprule
记录目录或文件 & 内容 \\\midrule
\texttt{output/results/deterministic} & 每个种子、条件数、预算与方法的逐步 CSV 及配置 JSON \\
\texttt{output/results/stochastic} & 噪声、批量、预算、真实梯度与动量误差 \\
\texttt{output/results/orthogonalization} & 各形状、输入谱、步数的误差与重复计时 \\
\texttt{output/results/language/pilot} & 全部候选学习率的试跑记录 \\
\texttt{output/results/language/formal} & 正式训练及测试指标，检查点步数和参数分组 \\
\texttt{data/manifest.json} & 固定数据链接、提交版本、SHA256、词表与划分大小 \\
\texttt{output/audit.json} & 运行计数、数值核查和结果记录集合校验值 \\
\bottomrule\end{tabularx}
\end{table}
正式检查点保存在本地输出目录，不纳入 GitHub 原始记录集合；复现脚本可以重新训练得到。计时相关字段会随机器负载变化，因而重新运行后的结果集合 SHA256 也会变化。校验值用于标识本次提交的记录集合，不要求不同机器重跑后具有相同的时间戳、耗时或集合摘要。

历史结果集合摘要为 \texttt{\ResultDigest}。附录语言模型使用的 Tiny Shakespeare 文件 SHA256 为
\begin{center}\footnotesize\ttfamily\expandafter\nolinkurl\expandafter{\DataSHA}\end{center}

\subsection{理论与实现的对照表}
\begin{table}[htbp]
\centering\small
\caption{理论版本与语言模型实现的对应关系}
\begin{tabularx}{\linewidth}{lXX}
\toprule 项目 & 理论与矩阵实验 & 语言模型实验 \\\midrule
正交化 & 精确紧致极因子 & 精确极因子或五步多项式 \\
动量初始化 & $M_1=G_1$ & 缓冲区初始化为零 \\
Nesterov / 形状缩放 & 均不使用 & 两种 Muon 配置均使用 \\
权重衰减 & 0 & 权重矩阵 0.01 \\
其他参数 & 单个矩阵 & AdamW 辅助参数组 \\
步长与动量 & 按定理设置 & 验证集选步长，$\beta=0.95$ \\
\bottomrule
\end{tabularx}
\end{table}
